% ============================================================================
% CHAPTER 6 — DISCUSSION
% Write after Chapter 5. Skeleton fixes the argument structure now.
% ============================================================================
\chapter{Discussion}
\label{ch:discussion}

\section{Interpreting the Ordering Effect}
\label{sec:disc-ordering}
% TODO: what the RQ1 Δ means; reconcile with Cheng'26 (their prompt-level
% TDD lost to Freeform; our gate-level TDD result may agree or disagree —
% either way the mechanism distinction carries the argument).

\section{The Verification Gap}
\label{sec:disc-verification}
% TODO: passing your own BRT is necessary, not sufficient (the L4-but-not-
% resolved cases); why a gold-free pipeline cannot fully close it (no oracle);
% what partially helps (breaker as measurement, EPR-style cross-validation as
% future work). Connect to Cheng'26 failure mode #3.

\section{Implications for Agent Design}
\label{sec:disc-design}
% TODO: architectural constraints vs prompt instructions; gate placement;
% what the negative results (critic/breaker fail-open) say about self-review
% without an oracle; capable-model requirement for autonomous loops.

\section{Threats to Validity}
\label{sec:disc-threats}
% TODO: internal — non-random development set (disclosed), single model,
% prompt sensitivity; external — Python-only, Lite's repo skew (django-heavy),
% benchmark contamination question; construct — resolved as ground truth,
% plausible-vs-correct.
